\medskip

\textbf{Partie A}

\medskip

On considère la fonction $f$ définie sur l'ensemble $]0~;~+ \infty[$ par

\[f(x) = 1 + x^2 - 2x^2\ln (x).\]

On admet que $f$ est dérivable sur l'intervalle et on note $f'$ sa fonction dérivée

\medskip

\begin{enumerate}
\item Justifier que $\displaystyle\lim_{x \to 0} f(x) = 1$ et, en remarquant que $f(x) = 1 + x^2(1 - 2\ln (x))$, justifier que $\displaystyle\lim_{x \to + \infty} f(x) = - \infty$.
\item Montrer que pour tout réel $x$  de l'intervalle $]0~;~+ \infty[$, \: $f'(x) = - 4x\ln (x)$.
\item Étudier le signe de $f'(x)$ sur l'intervalle $]0~;~+ \infty[$, puis dresser le tableau de variations de la fonction sur l'intervalle $]0~;~+ \infty[$.
\item Démontrer que l'équation $f(x) = 0$ admet une unique solution $\alpha$ dans l'intervalle $[1~;~+ \infty[$ et que $\alpha \in [1~;~\text{e}]$.
\end{enumerate}

On admet dans la suite de l'exercice, que l'équation $f(x) = 0$ n'admet pas de solution sur l'intervalle ]0~;~1].

\begin{enumerate}[resume]
\item On donne la fonction ci-dessous écrit en Python. L'instruction \emph{from lycee import *} permet d'accéder à la fonction ln.

\begin{center}
\begin{tabular}{>{\ttfamily}l}
from lycee import *\\
\\
def f(x) :\\
\qquad return 1 + x**2 - 2*x**2*ln(x)\\
\\
def dichotomie(p)\\
\qquad a=1\\
\qquad b=2.7\\
\qquad while b - a > 10**(-p) :\\
\qquad \begin{tabular}{|l}
\quad if f(a)*f((a+b)/2) < 0 :\\
\quad |\quad b = (a+b)/2\\
\quad else :\\
\quad |\quad a = (a+b)/2\\
\end{tabular}\\
\qquad  return (a,b)
\end{tabular}
\end{center}

\medskip

Il écrit dans la console d'exécution :

$>>>$ {\ttfamily dichotomie(1)}

Parmi les quatre propositions ci-dessous, recopier celle affichée par l'instruction précédente ? Justifier votre réponse (on pourra procéder par élimination).

\medskip

\begin{center}
\begin{tabular}{l l}
Proposition A :&(1.75, 1.9031250000000002)\\
Proposition B :&(1.85, 1.9031250000000002)\\
Proposition C :&(2.75, 2.9031250000000002)\\
Proposition D :&(2.85, 2.9031250000000002)\\
\end{tabular}
\end{center}
\end{enumerate}

\bigskip

\textbf{Partie B}

\medskip

On considère la fonction $g$ définie sur l'intervalle $]0~;~+ \infty[$, par 
\[g(x) = \dfrac{\ln (x)}{1 + x^2}.\]

On admet que $g$ est dérivable sur l'intervalle $]0~;~+ \infty[$ et on note $g'$ sa fonction dérivée.

On note $\mathcal{C}_g$ la courbe représentative de la fonction $g$ dans le plan rapporté à un repère \Oij.

\medskip

\begin{enumerate}
\item Démontrer que pour tout réel $x$ de l'intervalle $]0~;~+ \infty[$, \: $g'(x) = \dfrac{f(x)}{x\left(1 + x^2\right)^2}$.
\item Démontrer que la fonction $g$ admet un maximum en $x = \alpha$
\end{enumerate}

On admet que $g(\alpha) = \dfrac{1}{2\alpha^2}$.

\begin{enumerate}[resume]
\item On note $T_1$ la tangente à $\mathcal{C}_g$ au point d'abscisse 1 et on note $T_{\alpha}$ la tangente à $\mathcal{C}_g$ au point d'abscisse $\alpha$.

Déterminer, en fonction de $\alpha$, les coordonnées du point d'intersection des droites $T_1$ et $T_{\alpha}$.
\end{enumerate}

\bigskip
